\documentclass[11pt,a4paper]{article}

\usepackage[T1]{fontenc}
\usepackage{lmodern}
\usepackage[margin=28mm]{geometry}
\usepackage{amsmath,amssymb,amsthm}
\usepackage{booktabs}
\usepackage{microtype}
\usepackage[hidelinks]{hyperref}

\newtheorem{theorem}{Theorem}
\newtheorem{lemma}[theorem]{Lemma}
\newtheorem{corollary}[theorem]{Corollary}
\newtheorem{proposition}[theorem]{Proposition}
\theoremstyle{definition}
\newtheorem{keylemma}[theorem]{Key Lemma}
\theoremstyle{remark}
\newtheorem{remark}[theorem]{Remark}
\newcommand{\Ostar}{O^{*}}
\newcommand{\E}{\mathbb{E}}
\newcommand{\Prb}{\Pr}

\clubpenalty=10000
\widowpenalty=10000
\displaywidowpenalty=10000
\emergencystretch=2em
\hypersetup{pdftitle={A Simplified Plan for Max 2-SAT: One Search, One Count, One Key Lemma}}

\title{A Simplified Plan for Max 2-SAT:\\ One Search, One Count, One Key Lemma}
\author{}
\date{8 October 2026\\[3pt]\small Working notes}

\begin{document}
\maketitle

\begin{abstract}
We select one direction from the notes~\cite{N3,N4,N5,N6}: DPLL with the
pruning clauses of~\cite{N3} and the sum check of~\cite{N4}, along a fixed
variable order. Knuth's identity turns its running time into one probability,
that a uniformly random partial assignment survives $t$ steps. An exact
identity shows that the sum check is a race between the accumulated shortfall
of the clauses, a martingale, and the uncertainty left on the free variables.
Their expected values all scale like $\sqrt d$, which explains why the sum
check, unlike single clauses, cuts a typical partial assignment after about
$0.4n$ steps at every density, and why the search saves between $0.50n$ and
$0.62n$ in our runs up to $n=1600$. We state the one result we need, a Key
Lemma for random formulas of every density; the companion note~\cite{N8}
proves it with $\gamma\approx0.02$, for every graph with random signs. The
worst case reduces to the same argument under three cancellation conditions,
together with aggregation for formulas that violate them. Everything else is
parked.
\end{abstract}

\section{The choice}\label{sec:choice}

We keep one search: the pruning clauses of~\cite{N3} with the sum check
of~\cite{N4}. Three reasons.
\begin{enumerate}
 \item It is the only mechanism we have found that beats the timing
 barrier~\cite[Remark~29]{N3}: single clauses cut only near the end of their
 supports, while the sum decides early at every density (Tables~\ref{tab:death}
 and~\ref{tab:trees}).
 \item Its proofs are short. Soundness is already proved, and the running
 time is a single probability (Section~\ref{sec:count}).
 \item It isolates what remains. The sum fails exactly on aligned structure,
 where every rule implied by single flips fails~\cite[Section~9]{N4}, so the
 worst case needs one more ingredient, aggregation, and nothing else
 (Section~\ref{sec:worst}).
\end{enumerate}

\paragraph{The redundancy question of~\cite{N6}.} In this view a new clause
need not add information by itself. The sum adds up the shortfall of every
active clause, whether or not the clause is implied by the others. What makes
the sum decisive is cancellation: two clauses that contain opposite literals
of a shared variable make one of the two literals true whatever its value.
Only overlaps with the same literals defeat the sum, and if they dominate, the
formula is aligned.

\section{The search}\label{sec:search}

The formula is canonical with $B=1$, notation as in~\cite[Section~2]{N3}. Fix
the order $x_1,\dots,x_n$. A node at depth $t-1$ has $x_1,\dots,x_{t-1}$
assigned and branches on $x_t$, trying both values. Let $p$ be the smaller
deficit of $x_t$ at the node and $\ell$ the literal attaining it (ties: the
positive literal). The child in which $\ell$ is false receives the pruning
clause $P_t=\langle L_t\le p-1\rangle$ over the free neighbours of $x_t$; the
other child receives nothing. A node is cut when
\begin{itemize}
 \item some pruning clause has more true literals than it allows, or
 \item the \emph{sum check} fires: $\sum_{w\text{ free}}\min\{m^0_w,m^1_w\}
 >\sum_{c\in I}k_c$, where $I$ is the set of live pruning clauses that are not
 vacuous, $k_c$ is the allowance of $c$, and $m^b_w$ is the number of clauses
 of $I$ in which the value $b$ of $w$ makes a literal true~\cite[Section~5]{N4}.
\end{itemize}
There is no domination and no propagation; both are sound and help in
practice, but the analysis is for this plain version. If $p=0$, the pruned
child is cut at once. Soundness: each branching adds a clause to one child
only, so~\cite[Proposition~20]{N3} applies, and the sum is implied by the
clauses.

\section{Counting nodes}\label{sec:count}

\begin{proposition}[Knuth's identity~\cite{Knu75}]\label{prop:knuth}
Let $\sigma$ be uniform in $\{0,1\}^n$ and let $S_t$ be the probability that
the partial assignment $\sigma_1,\dots,\sigma_t$ survives all checks at depths
$1,\dots,t$, with $S_0=1$. The search tree has exactly $2^tS_t$ live nodes at
depth $t$, and at most $1+2n\max_t2^tS_t$ nodes in total. Hence the search runs
in time $\Ostar(\max_t2^tS_t)$ and polynomial space.
\end{proposition}

\begin{proof}
The checks at depth $s$ depend only on $\sigma_1,\dots,\sigma_s$, so the live
nodes at depth $t$ correspond to the surviving prefixes, each of probability
$2^{-t}$. Every node other than the root is a child of a live node, and each
live node has two children.
\end{proof}

So the whole running-time analysis is one probability estimate for a uniformly
random partial assignment. No measure, no clause weights, no tilted walk and no
availability argument are needed.

\section{The sum check as a race}\label{sec:race}

Let $c$ be the clause created at step $i$. Write $D_c$ for its support size,
the number of free neighbours of $x_i$ at that time, and $\Delta_c$ for the
effective gap of $x_i$ then, so that its threshold is $p_c=(D_c-|\Delta_c|)/2$
by~\cite[(1)]{N3}. At a later node, $a_c$ of its support variables are
assigned and $j_c$ of them took the value that makes their literal true. For a
set $I$ of clauses and a free variable $w$ put $b_w=m^1_w-m^0_w$.

\begin{lemma}[The race]\label{lem:race}
For every set $I$ of pruning clauses that are not violated,
\[
 \sum_{w\text{ free}}\min\{m^0_w,m^1_w\}-\sum_{c\in I}k_c
 =S(I)+Q(I)-U(I),
\]
where
\[
 S(I)=\sum_{c\in I}\Bigl(\frac{|\Delta_c|}2+1\Bigr),\qquad
 Q(I)=\sum_{c\in I}\Bigl(j_c-\frac{a_c}2\Bigr),\qquad
 U(I)=\frac12\sum_{w\text{ free}}|b_w| .
\]
So the sum check over $I$ fires if and only if $S(I)+Q(I)>U(I)$.
\end{lemma}

\begin{proof}
$\min\{m^0_w,m^1_w\}=\frac12(m^0_w+m^1_w)-\frac12|b_w|$, and
$\sum_w(m^0_w+m^1_w)=\sum_c(D_c-a_c)$. With
$k_c=p_c-1-j_c=\frac12(D_c-|\Delta_c|)-1-j_c$, each clause contributes
$\frac12(D_c-a_c)-k_c=\frac12|\Delta_c|+1+j_c-\frac12a_c$.
\end{proof}

$S$ is the \emph{shortfall}: every clause belongs to a variable set against
its majority, and its bound lies $|\Delta_c|/2+1$ below the expected count of
its support. $U$ is the \emph{uncertainty}: whatever values the free variables
take, the clauses that disagree about a variable cancel, and only $|b_w|$ is
left open. $Q$ records how far the assigned support variables deviated from
fair coins. The identity was checked in exact arithmetic on $7{,}259$ random
states.

\begin{remark}[Which clauses to sum]\label{rem:vacuous}
Removing a vacuous clause $v$ never weakens the check: its term
$\frac12(D_v-a_v)-k_v$ is at most $-\frac12(D_v-a_v)$, while removing it changes
$U$ by at most $\frac12(D_v-a_v)$. So the check of the search fires whenever
the check over all non-violated clauses fires. For the analysis we take $I_t$
to be all clauses created up to depth $t$ and write $S_t,Q_t,U_t$ and
$Z_t=S_t+Q_t-U_t$. A partial assignment that survives depth $t$ has
$Z_t\le0$.
\end{remark}

\begin{proposition}[One step]\label{prop:step}
Along a uniformly random $\sigma$, at step $t$:
\begin{enumerate}
 \item $Q_t-Q_{t-1}=\frac12(2\sigma_t-1)\,b_{x_t}$, with $b_{x_t}$ taken at
 depth $t-1$; so $Q$ is a martingale.
 \item $U$ loses $\frac12|b_{x_t}|$, because $x_t$ is no longer free.
 \item With probability $\frac12$, when $\sigma_t$ is the non-preferred value,
 a clause is created: $S$ grows by $\frac12|\Delta_t|+1$, and $b_w$ changes by
 $\pm1$ for each of its support variables $w$.
\end{enumerate}
In the random-sign model of Section~\ref{sec:key}, the signs added to the $b_w$
in item~3 are independent fair coins given the history, since the sign of $w$
in its clause with $x_t$ plays no role before step $t$. Hence
\[
 \E\bigl[Z_t-Z_{t-1}\bigm|\sigma_1,\dots,\sigma_{t-1}\bigr]
 =\frac{|\Delta_t|}4+\frac12+\frac{|b_{x_t}|}2-\frac{z_t}4,
\]
where $z_t$ is the number of free neighbours $w$ of $x_t$ with $b_w=0$.
\end{proposition}

\begin{proof}
Assigning $x_t$ raises $a_c$ for the clauses containing it and raises $j_c$ for
those in which $\sigma_t$ is the bad value, which gives item~1. Items~2 and~3
follow from the definitions. Adding an independent fair sign to $b_w$ raises
$|b_w|$ by one in expectation if $b_w=0$ and leaves it unchanged otherwise.
\end{proof}

\section{The Key Lemma}\label{sec:key}

\paragraph{Model.} In $R(n,d)$ each of the $\binom n2$ pairs carries one clause
with probability $d/(n-1)$, independently, with independent uniform signs. By
symmetry the fixed order does not matter.

\begin{keylemma}\label{kl}
There are constants $\gamma>0$ and $d_0$ such that for all $n$ and all
$d_0\le d\le n-1$, with probability $1-o(1)$ over $F\in R(n,d)$,
\[
 \max_{0\le t\le n}2^tS_t(F)\le2^{(1-\gamma)n}.
\]
\end{keylemma}

The Key Lemma is proved in~\cite[Corollary~3]{N8} for every $\gamma<0.0236$.
By Proposition~\ref{prop:knuth}, the search of Section~\ref{sec:search}
solves Max 2-SAT on $R(n,d)$ in time $\Ostar(2^{(1-\gamma)n})$ and polynomial
space, with probability $1-o(1)$, for every density. This is an average-case
statement, and its novelty still has to be checked against the literature.

\paragraph{Evidence.} Table~\ref{tab:death} gives the depth at which a
uniformly random partial assignment is cut, Table~\ref{tab:trees} the size of
the whole tree, both in exactly this setting.

\begin{table}[ht]
\centering
\begin{tabular}{llcccc}
\toprule
 & & $d=16$ & $64$ & $256$ & $(n-1)/2$\\
\midrule
$n=400$ & single clauses & 0.19 & 0.52 & 0.76 & 0.72\\
 & with the sum check & 0.19 & 0.41 & 0.43 & 0.43\\
$n=1600$ & single clauses & 0.11 & 0.49 & 0.72 & 0.84\\
 & with the sum check & 0.11 & 0.39 & 0.42 & 0.44\\
\bottomrule
\end{tabular}
\caption{Median depth, as a fraction of $n$, at which a uniformly random
partial assignment is cut; $30$ assignments per entry.}\label{tab:death}
\end{table}

\begin{table}[ht]
\centering
\begin{tabular}{llccc}
\toprule
 & & $d=64$ & $256$ & $(n-1)/2$\\
\midrule
$n=400$ & single clauses & 0.36 & 0.20 & 0.26\\
 & with the sum check & 0.52 & 0.51 & 0.54\\
$n=800$ & single clauses & 0.40 & 0.24 & 0.21\\
 & with the sum check & 0.56 & 0.53 & 0.50\\
$n=1600$ & single clauses & 0.43 & 0.24 & 0.16\\
 & with the sum check & 0.62 & 0.56 & 0.56\\
\bottomrule
\end{tabular}
\caption{Saving $n-\log_2(\text{tree size})$ as a fraction of $n$, estimated by
population splitting (unbiased; one formula per entry).}\label{tab:trees}
\end{table}

With single clauses the cut moves towards the last level as $d$ grows; with the
sum check it stays near $0.4n$, and the saving stays between $0.50n$ and $0.62n$.

\paragraph{Why the density cancels.} For $d\to\infty$ use the approximations
that a free variable at time $\tau n$ lies in about $\tau d/2$ clauses, that
$\Delta_t$ has variance about $d(1+\tau)$, and that $b_w$ is a sum of
independent signs. Then Proposition~\ref{prop:step} gives drift $\sqrt d\,\mu(\tau)$
per step with
\[
 \mu(\tau)=\frac14\sqrt{\frac{2(1+\tau)}\pi}+\frac12\sqrt{\frac\tau\pi}
 -\frac{1-\tau}{4\sqrt{\pi\tau}},
\]
and $Z_{\tau n}\approx n\sqrt d\,F(\tau)$ with
\[
 F(\tau)=\int_0^\tau\mu
 =\frac1{\sqrt\pi}\Bigl(\frac{\sqrt2}6\bigl((1+\tau)^{3/2}-1\bigr)
 +\frac{\tau^{3/2}-\tau^{1/2}}2\Bigr).
\]
Early on, new clauses meet fresh variables and $U$ grows faster than $S$; the
drift turns positive at $\tau_0\approx0.14$, and $F$ changes sign at
$\tau^*\approx0.46$, independently of $d$. The measured cuts at $0.39$--$0.44$
come slightly earlier, because single clauses and the removal of vacuous
clauses also act. A single clause, by contrast, compares an imbalance of order
$\sqrt d$ with a remaining support of order $d$, and so decides only at depth
$1-O(1/\sqrt d)$.

\section{How it is proved}\label{sec:proof}

The proof in~\cite{N8} is simpler than the outline we had here (drift, noise,
combine).
\begin{enumerate}
 \item \textbf{One depth.} $\max_t2^tS_t\le\max\{2^T,2^nS_T\}$, so it
 suffices to show that a random partial assignment survives one depth
 $T=\tau n$ only with probability $e^{-cn}$; then
 $\gamma=\min\{1-\tau,\,c/\ln2\}$.
 \item \textbf{Late signs.} The search never uses the signs of $x_t$ before
 it reaches $x_t$. So, given the past, $\Delta_t$ is a sum of unseen signs
 ($\pm1$ over the 2-clauses of $x_t$, $\pm2$ over its unit clauses), the bad
 values of $x_t$ in the clauses that contain it are unseen fair coins, and at
 depth $T$ so are the bad values of the free variables.
 \item \textbf{Three tails.} Szarek's inequality and a one-sided bound for
 nonnegative increments give the shortfall; $Q$ is a walk with fair steps
 $\pm\frac12$ (Azuma); $U$ is controlled by Jensen's and McDiarmid's
 inequalities. No drift analysis of $U$ is needed, because $U$ is evaluated
 only at depth $T$, where the bad values of the free variables are still
 unseen.
\end{enumerate}
The graph enters only through a few degree sums, so the result holds for every
graph with random signs whose margin is of order $n\sqrt d$: every
$\gamma<0.0236$ for $R(n,d)$, and every $\gamma<0.0132$ for every $d$-regular
graph in any order~\cite[Corollaries~3 and~4]{N8}. The constant is far below
the measured $0.50n$--$0.62n$ and the Gaussian estimate of about $0.2n$ from the
drift above. The loss comes from four separately added tails, from Szarek's
constant, and from using a single depth. A single exponential moment of $Z_t$
near depth $0.5n$ is the natural way to improve it.

\section{The worst case}\label{sec:worst}

The proof of~\cite{N8} uses random signs in three places: the size and
concentration of the shortfall, the step size of the martingale $Q$, and the
size of the uncertainty $U$. For a fixed formula, $\sigma$ alone keeps $Q$ a
martingale and gives the shortfall its mean, but nothing more. In particular
our earlier claim that the shortfall part holds for every formula is true only
for its mean.

\paragraph{Cancellation conditions~\cite[Section~6]{N8}.} The proof goes
through verbatim for a fixed formula if, except with probability
$e^{-\Omega(n)}$ over $\sigma$: (W1) $U_T$ is a small constant times
$n\sqrt d$; (W2) $\sum_{t\le T}b_{x_t}^2=O(nd)$; (W3) the effective gaps keep
their typical size, $\sum_t|\Delta_t|=\Omega(n\sqrt d)$ and
$\sum_t\Delta_t^2=O(nd)$. (W1) and (W2) say that the clauses containing a
variable disagree about its value often enough that $|b_w|$ is of order
$\sqrt{m_w}$ rather than $m_w$, at the free and at the assigned variables.

\paragraph{When it fails.} If cancellation fails, many clauses agree about the
values of many variables: the formula is aligned, as in the circulant formulas
$C_{n,c}$, where every clause has the same bad value at every variable. By the
barrier of~\cite[Section~9]{N4}, no rule implied by single flips can help
there: Max-Cut on $K_m$ has $\binom m{m/2}$ strict local maxima, and $C_{n,c}$
has exponentially many optimal assignments; even comparisons in both children
with fixed tie-breaking keep $126$ assignments of $C_{9,2}$, more than the $70$
leaves of plain domination. Both are solved by aggregation: counting the ones
within exact interaction classes~\cite[Section~5]{N5}.

\paragraph{The worst-case target.} (i) The Key Lemma under (W1)--(W3); this
is the remark above. (ii) Formulas that violate them have enough aligned
structure for aggregation to pay. (iii) A rule that decides between the two at
a node. Part (ii) is the hard part; it is a statement about formulas whose
active clauses agree about their variables, not about arbitrary graphs. By
Corollary~4 of~\cite{N8} the graph alone is never the obstacle.

\section{Parked}\label{sec:parked}

Clause-weight measures and the availability analysis~\cite[Sections~6--8]{N4};
packing and core certificates, coverage, exposure, the entropy dichotomy and the
desired frontier theorem~\cite{N5,N6}; comparisons in both children~\cite{N6};
pair decisions~\cite[Section~3]{N4} except as a possible form of aggregation.
None of them was needed for the Key Lemma. The barrier and the
linear-relaxation characterization of sums~\cite{N4} are kept, because they
define the worst case and the strongest version of the check.

\section{Checks}\label{sec:checks}

\begin{itemize}
 \item \texttt{explore\_death\_depth.py} (about $15$ seconds):
 Table~\ref{tab:death}.
 \item \texttt{explore\_key\_lemma.py}: \texttt{check} verifies
 Lemma~\ref{lem:race} in exact arithmetic; \texttt{400 50} and similar give
 the rows of Table~\ref{tab:trees} (minutes for $n\le800$, longer for $1600$).
 \item \texttt{explore\_key\_lemma\_proof.py}: the checks of~\cite{N8}.
\end{itemize}

\begin{thebibliography}{9}
\raggedright
\bibitem{N3} Max 2-SAT and the domination rule: results and ideas so far.
 Working note, 5 October 2026 (\texttt{max2sat\_results\_and\_ideas.tex}).
\bibitem{N4} Max 2-SAT with pruning clauses: pair decisions, sums, and a
 measure. Working note, 6--7 October 2026
 (\texttt{max2sat\_pairs\_sums\_measure.tex}).
\bibitem{N5} Exploring the branching frontier for Max 2-SAT. Draft,
 7 October 2026 (\texttt{max2sat\_branching\_frontier.tex}).
\bibitem{N6} A learning and coverage dichotomy for Max 2-SAT. Draft,
 7 October 2026 (\texttt{max2sat\_learning\_dichotomy.tex}).
\bibitem{N8} The Key Lemma for random signs. Working note, 8 October 2026
 (\texttt{max2sat\_key\_lemma\_proof.tex}).
\bibitem{Knu75} D.~E.~Knuth. Estimating the efficiency of backtrack
 programs. \emph{Mathematics of Computation} 29(129):121--136, 1975.
\end{thebibliography}

\end{document}
